\scanpage{45}
% archive sheet number: 44

(où $\unsure{C'|\,'}$ signifie exact à g) et $\Lambda$-alg. à
$\underline{\operatorname{Hom}}'(R_\Lambda,\mathrm{Ens})$ où $R_\Lambda$ cat. des schémas de t.f. sur $\Lambda$, le foncteur d'inclusion

\diagram{\resizebox{\linewidth}{!}{\begin{tikzcd}[ampersand replacement=\&,column sep=small,row sep=small]
 \&\underline{\operatorname{Hom}}'(R_\Lambda,\mathrm{Ens})\arrow[r,hook]\&
 \underline{\operatorname{Hom}}'(R,\mathrm{Ens},\,)\\
 {A\mapsto(B\mapsto\operatorname{Hom}_{\Lambda}(B,A))}\&
 (\Lambda\text{-alg})\arrow[u,"\sim"]\&
 (\mathrm{Ann})\arrow[u,"\sim"]\arrow[l,hook]\&
 {A\mapsto(B\mapsto\operatorname{Hom}_{\mathbb Z}(B,A))}\\
 \Lambda\text{-alg}\&\&\mathbb Z\text{-alg}\&\shortstack{$\mathbb Z$-alg.\\$\parallel$\\$\operatorname{Hom}_\Lambda(B_\Lambda,A)$}
\end{tikzcd}}}

est transposé du foncteur $R\longrightarrow R_\Lambda$, $X\to X\otimes_{\unsure{e|\mathbb Z}}\Lambda$ i.e.
$X\mapsto X\times_e S=X\times S$. \inserted{si $A$ un $\Lambda$-alg.}
(qui est en effet exact à g). Donc

\[
 X_{\Lambda\text{-alg}}\simeq\widetilde R_{\Lambda\mathrm{top}}\simeq\widehat R_\Lambda
\]

$\Lambda$-top : $X_i\to X$ couvrante ssi les $X_i\deleted{\unsure{\otimes|\times}}\times_e S\to X\times_e S$ ont une section.
\uline{NB} $R_\Lambda$ est aussi déduit de $R$ par cat. des fractions, en inversant les $f$ telles que $f\times_e S$ un isom.

\[
\begin{aligned}
 \text{\uline{Ex}}\quad &X_{\mathbb Z/n\mathbb Z\text{-alg}},\quad X_{\mathrm{car}\ p}\quad
 (\text{cas particulier}),\quad X_{\mathbb Z[1/n]}\\
 &X_{\mathbb Q}=X_{\mathrm{car.}\,0}=\deleted{\unsure{\mathrm{car}\neq p}}\bigcap_{n\geq2}X_{\mathbb Z[1/n]}
\end{aligned}
\]

\uline{NB} Ce qu'on vient de voir se généralise au cas d'une $R$ quelconque avec $\varprojlim$ finies, et l'espèce de structure correspondante : $T$.
On se donne $\Lambda\in T(\mathrm{Ens})\simeq\operatorname{Ind}R^\circ$, i.e. $S\in\operatorname{Pro} R$,
on suppose que $S\to e_R$ \deleted{\unsure{un}} mono, et on considère $R_S=$ catégorie des fractions de $R$ par les flèches qui deviennent des isom par $\times_e S$,
donc sous-catégorie image ess. de $\operatorname{Pro}R_{/S}$ par $R\to\operatorname{Pro}R_{/S}$, $X\mapsto X\times_e S$. \deleted{\unsure{par}}.
Comme $\operatorname{Pro} R\hookrightarrow\widehat R$, \deleted{\unsure{on peut}} on peut supposer plus gén. $S\in\operatorname{Ob}\widehat R$, \inserted{$S\subset e_{\widehat R}$}
pas néc. dans l'image essentielle de $\operatorname{Pro} R$, et s'intéresser aux $f:\mathcal E\to\widehat R$ tels que
$f^*(\deleted{\unsure{A}}S)\to f^*(e_{\widehat R})=e_{\mathcal E}$ se factorise par $f^*(S)\hookleftarrow e$.
C'est la forme générale des ss-structures qui sont représentées par des \uline{ouverts} du topos classifiant $\widehat R$.

% slanted
\margin{$T(\mathrm{Ens})^\circ$ \inserted{exact à g.}\quad montrer que si $S\in\operatorname{Pro}R$, \unsure{$\exists$ là,} c'est la condition précédente}
